%%%PAGE 113 rot=0 legibility=good summary=电磁波谱与光谱分类、折射与全反射的应用举例（彩虹、视深、光斑环面积、光导纤维、水槽光线与水位上升速度）
%%%BLOCK id=113-01 ch=14 sec=电磁波谱 kind=summary alt=none cont=none y=0.04-0.14
无线电波\quad 微波\quad 红外线\quad 可见光\quad 紫外线\quad $X$射线\quad $\gamma$射线 $\longrightarrow$\quad $f\uparrow$\quad $\lambda\downarrow$

\fbox{无微红可紫 $X$ $\gamma$}
%%%END
%%%BLOCK id=113-02 ch=17 sec=光谱 kind=summary alt=14,15 cont=none y=0.14-0.21
% CHECK: 最左侧竖排三字疑为“连续谱”，其后大括号内为“吸收谱/发射谱”，层级关系按原稿位置转录
连续谱\ $\left\{\begin{tabular}{@{}l@{}}吸收谱\\发射谱\end{tabular}\right.$\ 反映内部结构光谱\ $\left\{\begin{tabular}{@{}l@{}}连续谱\\离散谱\end{tabular}\right.$
%%%END
%%%BLOCK id=113-03 ch=15 sec=折射·全反射 kind=knowledge alt=none cont=none y=0.22-0.35
\begin{enumerate}
\item 彩虹为光的色散/折射
\item \unclear{鸟}看鱼偏高、鱼看鸟偏高
\item 水中唱歌\quad 从空气到水中\ 声波波长$\uparrow$，电磁波波长短
\item 海\unclear{市蜃}楼、车尾灯光带(全反射)
\end{enumerate}
%%%END
%%%BLOCK id=113-04 ch=15 sec=折射·全反射 kind=problem alt=none cont=none y=0.35-0.45
% 原稿第5题：只有图与解答，没有写出题干
\begin{problem}[5]
%FIG p113_a 0.09 0.352 0.205 0.43
\notefig[0.25]{figures/p113_a.png}
\begin{solution}
\begin{align*}
&\sin\alpha=\frac{1}{n_1}=\frac{r_1}{\sqrt{r_1^2+d^2}}\quad\therefore\ r_1=\frac{d}{\sqrt{n_1^2-1}}\qquad
\sin\beta=\frac{1}{n_2}=\frac{r_2}{\sqrt{r_2^2+d^2}}\qquad r_2=\frac{d}{\sqrt{n_2^2-1}}\\
&S=\pi r_2^2-\pi r_1^2=\pi d^2\left(\frac{1}{n_2^2-1}-\frac{1}{n_1^2-1}\right)
\end{align*}
\end{solution}
\end{problem}
%%%END
%%%BLOCK id=113-05 ch=15 sec=折射·全反射 kind=derivation alt=none cont=none y=0.45-0.66
\textbf{6．光导纤维．}

%FIG p113_b 0.055 0.478 0.25 0.567
\notefig[0.35]{figures/p113_b.png}

$n_1>n_2$
\[n_1\sin\theta_3=n_2\sin\frac{\pi}{2}\qquad\therefore\ \sin\theta_3=\cos\theta_2\qquad\therefore\ \theta_3+\theta_2=\frac{\pi}{2}\]
\[\left\{\begin{array}{l}n_1=\dfrac{\sin\theta_1}{\sin\theta_2}\\[2mm]\sin^2\theta_2+\cos^2\theta_2=1\quad\therefore\ \sin\theta_1=\sqrt{n_1^2-n_2^2}\end{array}\right.\]

路程\ $x=\dfrac{L}{\cos\theta_2}$\qquad $v=\dfrac{c}{n_1}$\qquad $t=\dfrac{x}{v}=\dfrac{n_1^2}{n_2^2}\,\dfrac{L}{c}$ % CHECK: 分母似为 n_2^2，但由前式 cosθ2=n2/n1 推得应为 n_1^2/n_2 · L/c

% 以下两式在图的左下方；第一式整体被划去
\[\text{\struck{$\dfrac{1}{\sin\theta_3}=n=\dfrac{\sin\theta_1}{\sin\theta_2}$}}\]
\[\frac{n_1}{n_2}=\frac{1}{\sin\theta_3}\]
%%%END
%%%BLOCK id=113-06 ch=15 sec=折射·全反射 kind=problem alt=none cont=none y=0.66-0.92
% 原稿第7题：只有图与解答，没有写出题干
\begin{problem}[7]
% 图p113_c内含标注：中点、中线、H、ΔH、V0、水位上升速度V0、D（底面宽）、Δd、d、α、β
%FIG p113_c 0.0 0.675 0.262 0.832
\notefig[0.4]{figures/p113_c.png}
\begin{solution}
① 求折射率
\begin{align*}
&\sin\alpha=\frac{D}{\sqrt{D^2+H^2}}\qquad \sin\beta=\frac{d-\frac{D}{2}}{\sqrt{\left(d-\frac{D}{2}\right)^2+\left(\frac{H}{2}\right)^2}}\\
&n=\frac{\sin\alpha}{\sin\beta}=\frac{D}{2d-D}\sqrt{\frac{(2d-D)^2+H^2}{D^2+H^2}}
\end{align*}
② 求光源与水面速度关系
% 下式（①式）下方有红色波浪线
\begin{align*}
&\Delta d=\Delta H(\tan\alpha-\tan\beta)\quad\text{①}\\
&\tan\alpha=\frac{D}{H}\qquad \tan\beta=\frac{d-\frac{D}{2}}{\frac{H}{2}}\quad\nearrow\\
&\Delta d=\Delta H\,\frac{(2H-2d)}{H}\quad\text{②}\\
&\frac{\Delta d}{\Delta t}=\frac{\Delta H}{\Delta t}\,\frac{(2\unclear{H}-2d)}{H}\quad\text{③}\\
&\therefore V=\frac{2(D-d)}{H}V_0\quad\text{④}
\end{align*}
% CHECK: ②③式分子中的“2H”疑应为“2D”（由 tanα-tanβ=(2D-2d)/H 及④式2(D-d)/H）；③式中该字母似为D写在H之上。tanβ 行末有一弯曲箭头指向上方①式
\end{solution}
\end{problem}
%%%END
%%%PAGEEND 113
